% Regression: a string bends at its waypoints, so it leaves each end toward
% the waypoint on that end's own side of the route and not toward the far
% end the chord would suggest.  A station read from the far end reserves a
% face the ink never uses and leaves the one it does use looking free.
\documentclass{standalone}
\usepackage{tenkz}
\makeatletter
\ExplSyntaxOn
\file_input:n { tenkz-kernel.code.tex }
\int_new:N \l__tenkz_test_side_count_int
\tl_new:N \l__tenkz_test_side_label_tl
\cs_new_eq:NN
  \__tenkz_test_side_auto:nN \__tenkz_kernel_r_label_auto_side:nN
\cs_set_protected:Npn \__tenkz_kernel_r_label_auto_side:nN #1#2
  {
    \__tenkz_test_side_auto:nN {#1} #2
    \int_gincr:N \l__tenkz_test_side_count_int
    \__tenkz_model_get:nnN {#1} {label} \l__tenkz_test_side_label_tl
    \str_case:VnF \l__tenkz_test_side_label_tl
      {
        {V}
          {
            \str_if_eq:VnF #2 {n}
              { \errmessage{a~routed~string~left~the~face~it~turns~into~free} }
          }
        {Y}
          {
            \str_if_eq:VnF #2 {n}
              { \errmessage{a~route~lost~the~face~its~first~waypoint~takes} }
          }
        {Z}
          {
            \str_if_eq:VnF #2 {n}
              { \errmessage{a~route~did~not~read~its~last~waypoint} }
          }
        {X}
          {
            \str_if_eq:VnF #2 {s}
              { \errmessage{a~straight~wire~read~a~waypoint~it~never~draws} }
          }
      }
      { \errmessage{unexpected~label~in~routed-wire~fixture} }
  }
\ExplSyntaxOff
\makeatother
\begin{document}
\tenkzkernel
% One waypoint: the route leaves V southward, so V gives up its south.
\begin{tenkz}[lattice={2x3}, bonds=none]
  \tn[at=(1,1), skin=dot]{V}
  \tn[at=(1,3)]{}
  \tn[at=(2,1), skin=none, name=jV1]{}
  \tnwire[kind=string, name=segV1]{(1,1)}{jV1}
  \tnwire[kind=string, name=segV2,
          cross={over at crossing of segV1 and segV2}]{jV1}{(1,3)}
\end{tenkz}
% Two waypoints: each end turns into the one below it, so both ends give up
% their south.  Reading the first waypoint at both ends would leave Z's.
\begin{tenkz}[lattice={2x4}, bonds=none]
  \tn[at=(1,1), skin=dot]{Y}
  \tn[at=(1,4), skin=dot]{Z}
  \tn[at=(2,1), skin=none, name=jYZ1]{}
  \tn[at=(2,4), skin=none, name=jYZ2]{}
  \tnwire[kind=string, name=segYZ1]{(1,1)}{jYZ1}
  \tnwire[kind=string, name=segYZ2,
          cross={over at crossing of segYZ1 and segYZ2}]{jYZ1}{jYZ2}
  \tnwire[kind=string, name=segYZ3,
          cross={over at crossing of segYZ2 and segYZ3}]{jYZ2}{(1,4)}
\end{tenkz}
% An index wire draws the chord between its ends, so its ends read each
% other and X keeps its free south.
\begin{tenkz}[lattice={2x3}, bonds=none]
  \tn[at=(1,1)]{}
  \tn[at=(1,3), skin=dot]{X}
  \tnwire{(1,1)}{(1,3)}
\end{tenkz}
\ExplSyntaxOn
\int_compare:nNnF { \l__tenkz_test_side_count_int } = {4}
  { \errmessage{routed-wire~fixture~did~not~station~four~labels} }
\ExplSyntaxOff
\end{document}
